\uuid{mRoq}
\titre{Implication/équivalence}
\auteur{bodin ; hanani ; mzari}
\organisation{exo7}
\licence{CC BY-NC-SA 4.0}
\datecreate{2020-01-20}
\niveau{L1}
\module{Algèbre}
\chapitre{Logique, ensemble, raisonnement}
\sousChapitre{Logique}
\difficulte{3}
\theme{Logique, implication, équivalence}
\exo7qcmid{lille-1.logique.015}

\contenu{
\begin{question}
\qclassification{100.01}
\qtags{motcle=implication/équivalence}
À quoi est équivalent $P \Longleftarrow Q$ ?
\begin{answers}
    \good{non($Q$) ou $P$}

    \bad{non($Q$) et $P$}

    \bad{non($P$) ou $Q$}

    \bad{non($P$) et $Q$}
\end{answers}
\begin{explanations}
La définition (à connaître) de "$P \implies Q$" est "non($P$) ou $Q$".
\end{explanations}
\end{question}
}
